\documentclass[stu,12pt,floatsintext]{apa7}

\usepackage[american]{babel}
\usepackage{csquotes}
\usepackage[style=apa,sortcites=true,sorting=nyt,backend=biber]{biblatex}
\DeclareLanguageMapping{american}{american-apa}
\addbibresource{bibliography.bib}

\usepackage[T1]{fontenc}
\usepackage{mathptmx}
\usepackage{booktabs}
\usepackage{url}

\title{Coal-to-X House-hold Level Clean Heating Transition Simulation Sandbox}
\shorttitle{Coal-to-X Simulator}
\author{Guo Hang}
\duedate{[TODO: submission date]}
\affiliation{PKUHS Stanford}
\course{CTB 2025--26}
\professor{[TODO: instructor name]}

\abstract{%
Rural North China's coal-to-X transition has reduced coal pollution, but many households still struggle with winter heating costs as subsidies taper. This study asked whether a household-level interactive simulation could help users understand trade-offs among clean air, affordability, and compliance, while producing evidence for village pathway choice. We built and deployed a five-turn clean-heating sandbox in which users entered household income, housing area, and heating costs, compared gas and heat-pump routes, and viewed energy burden, compliance, and emission indicators. From June 28 to July 5, 2026, 21 participants completed the simulation and post-survey; 14 non-farmer users also had paired pre/post understanding scores. Understanding improved for 11 of those 14 users, 17 participants rated the AI debrief as helpful, and farmer logs showed that route choices and outcomes differed by household conditions. The sandbox made policy trade-offs concrete and suggests that village clean-heating decisions should compare household-level affordability rather than assume one universal route.%
}

\begin{document}
\maketitle

\section{Introduction}

Winter heating is a basic livelihood need in rural North China. Since the national push for clean winter heating in the mid-2010s, tens of millions of farm households have shifted from scattered coal to coal-to-X routes such as natural gas, electricity, or heat pumps. Regional air quality has improved and winter coal use has fallen sharply \parencite{ndrc2017plan,yuan2025carbon}. Environmental gains, however, do not automatically translate into sustainable energy use for every family. Upfront retrofit costs and recurring winter bills are largely borne by households; as equipment and operating subsidies phase down by cohort, many face a familiar tension: households can afford the switch but not everyday use. Heating bills rise as a share of annual income, indoor temperatures drop, and some households return to coal or maintain dual fuel sources \parencite{he2021decoalizing,zhai2023review}. Public debate often still asks only whether clean heating is ``good,'' rather than which coal-to-X pathway fits diverse household budgets, housing conditions, and compliance constraints.

The policy stakes are concrete. In much of North China, village- or township-level campaigns coordinate technology choice and installation. If the chosen route does not match most households' cash flow and energy burden, transitions that look successful on paper may prove unstable in operation. Prior work documents this tension from multiple angles. Environmental and health studies show that clean heating reduces particulate exposure and carbon emissions. Energy-economics and energy-poverty research finds that gas spending above roughly 5\% of disposable income materially changes behavior, and that subsidy withdrawal raises winter costs and the share of households under affordability stress \parencite{zhao2024sustainable}. Regional cost--benefit analyses further show that health gains and economic costs are unevenly distributed across cities \parencite{zhang2024uneven}. Agent-based and regional policy models link household actions to village- or city-scale outcomes \parencite{wang2023abm,liu2020costs}, yet few public-facing tools let users enter household income, floor area, and heating costs, compare gas, heat-pump, and insulation routes interactively, and experience subsidy phase-out and compliance pressure in one session.

This study addresses that gap by designing, deploying, and evaluating a \textit{village-level clean-heating transition simulation sandbox} (\url{https://www.clean-heating-simulator.com}). Participants role-play a farm household, enter whole-household income, surplus, floor area, and pre-transition heating cost, then make choices across five turns---whether to transition, which technology to adopt, how to boost income or save energy, and how to cope with subsidy withdrawal---while viewing live indicators for heating cost, energy burden rate, legal compliance, and emission score. After the terminal state, they may request an AI debrief generated from session logs via a live language-model API; no curated mock script is served if the API fails. The simulation engine is deterministic and calibrated from literature and public statistics.

Our research question is: \textit{Can household-level interactive simulation help users with different identities---students, converted and unconverted farmers, and other stakeholders---understand trade-offs among environmental goals, heating affordability, and compliance, and supply aggregatable evidence for village pathway comparison?} We hypothesize that (H1) farmer users can identify coal-to-X paths better matched to their own household conditions, and (H2) students and other non-farmers show improved understanding of heating-cost dynamics after one session. From June 28 to July 5, 2026, we recruited 21 fully completed cases (including rural users from Hebei Province and Beijing municipalities) through course outreach and a public URL, using pre/post surveys and behavioral logs.

The remainder of the paper is organized as follows. The Method section describes the deployed artifact, design, consent, and analysis plan. Results and Discussion sections will report survey, log, and simulated-outcome findings and their implications for village clean-heating decisions.

\section{Method}

\subsection{Participants}

We targeted 20+ participants, including classmates and rural users [TODO: confirm village-cadre-assisted recruitment], plus organic traffic via the public deployment URL. [TODO: recruitment period, final $N$, target village name.] Simulator constants were additionally calibrated from prior field conversations in rural North China and from published literature [TODO: interview $N$ and whether reported separately].

\subsection{Materials}

We deployed a village-level clean-heating transition simulation sandbox for rural North China coal-to-X contexts (\url{https://www.clean-heating-simulator.com}; single-page web application hosted on Vercel). Participants needed no account. After selecting an identity on the welcome screen (student, converted farmer, unconverted farmer, or other stakeholder) and entering basic background, they created a North China farm household profile: whole-household annual income, annual surplus, floor area, resident population, and pre-transition winter heating cost. Role (farm household), region (North China), and initial heating mode (scattered coal) were fixed.

The simulation ran five turns: (1)~status quo choice (continue coal or prepare transition); (2)~compliance pressure (probabilistic home visit if coal continues, scaled by floor area and consecutive non-compliance); (3)~technology route (natural gas, ground-source heat pump, or air-source heat pump; retrofit and operating costs linearly scaled from a 100\,m\textsuperscript{2} baseline); (4)~income boost or energy saving (one of three options); (5)~subsidy phase-out (equivalent gas volume derived from initial coal spending, then a fixed per-m\textsuperscript{3} cost increase, followed by terminal adjudication). The interface displayed six live indicators---annual heating cost, income, surplus, legal compliance (0--100), emission score (0--100), and energy burden rate (\%)---plus a turn-by-turn narrative log. Terminal outcomes were successful transition, financial strain, or enforcement-related seizure.

After the terminal state, participants could request an AI debrief: the browser submitted initial inputs, stepwise logs, and final state to a server endpoint; the server called the DeepSeek Chat Completions API (\texttt{deepseek-v4-flash}) with one of three human-authored structured system prompts (converted-farmer comparison, unconverted-farmer retrofit advice, or default analysis). The model was instructed to use only logged numbers and return three Markdown sections. If the API key was missing or the call failed, the interface showed an error---no curated mock script was served as fallback. The simulation engine itself was deterministic; economic, compliance, and emission constants were literature-calibrated, with some thresholds simplified as noted in-app.

Each session was stored anonymously (UUID) in a Supabase database table (\texttt{simulation\_sessions}), including baseline inputs, turn-choice IDs and labels, event logs, terminal type and summary, final metrics, AI analysis text (if generated), and post-survey fields. Table~\ref{tab:material-boundaries} summarizes component boundaries.

\begin{table}[htbp]
    \centering
    \caption{Material Boundaries: Simulation Engine, AI Debrief, and Empirical Logging}
    \label{tab:material-boundaries}
    \begin{tabular}{@{}p{0.22\linewidth}p{0.28\linewidth}p{0.42\linewidth}@{}}
        \toprule
        Component & Nature & Note \\
        \midrule
        Five-turn engine & Deterministic algorithm with literature constants & Outcomes computed by rules and probability models, not by an LLM \\
        AI debrief & Live DeepSeek API & Human-curated system prompts; output varies by session log; no mock fallback on failure \\
        Questionnaires and logs & Empirical capture & Stored in Supabase for pre/post and pathway analysis \\
        \bottomrule
    \end{tabular}
\end{table}

\subsection{Measures}

Dependent variables fell into three layers. \textit{Survey cognition and attitudes:} for students and other non-farmers, pre Q1 and post Q2 (5-point Likert on understanding of coal-to-gas and heating costs), Q4 (strategy change if replayed), S1 (simulator helpfulness vs.\ text), recommendation intent, and AI helpfulness; for converted farmers, cost realism (F4), best- and worst-matching stages (F6--F7), and subsidy coping; for unconverted farmers, concerns (N2), attitude shift (N3), and preferred route (N4). \textit{Behavioral logs:} five-turn choices and terminal type. \textit{Simulated outcomes:} terminal energy burden, heating cost, surplus, compliance, emission score, and CO\textsubscript{2} tons.

Open-ended post-survey text (e.g., student reflection, policy conflict detail, improvement suggestions, and negative AI feedback) will be analyzed via [TODO: coding scheme, e.g., dual-coder thematic analysis / codebook]; no codebook is fixed yet.

\subsection{Design}

We used a single-session exposure design combining pre/post questionnaires and behavioral logging. The intervention was completing one full sandbox run (optional AI debrief and post-survey).

Independent variables were (a)~exposure to the live sandbox and self-selected pathway (technology, compliance branch, income/energy choices---not experimenter-assigned) and (b)~participant identity (student, converted farmer, unconverted farmer, other) as a between-subjects factor.

Non-farmers completed a one-item pre-survey (Q1) before simulation and a approximately 2-minute post-survey after AI debrief. Converted and unconverted farmers skipped the pre-survey and only completed the post-survey (optional interview willingness; phone number collected only if interview was accepted).

Comparison targets included within-person pre/post deltas (Q2$-$Q1, non-farmers only); distributions of burden, compliance, and emission by route and terminal type; and cross-household comparison within the same village (\texttt{village\_name}) [TODO: whether village-level aggregation is implemented]. Exported Supabase records will be analyzed with [TODO: paired pre/post test, e.g., Wilcoxon signed-rank or paired $t$-test with normality check] for Likert items; categorical outcomes (route, terminal type, attitude shift) via [TODO: descriptive frequencies / Fisher's exact / $\chi^2$]; and simulated metrics compared across routes and identities [TODO: specific grouping and tests]. Thematic coding will supplement quantitative findings [TODO: whether themes enter formal inferential models].

\subsection{Procedure}

Participation began with three first-screen layers: welcome (identity and background), context briefing, and informed consent. The consent screen stated that data were anonymous (no name; phone not collected by default), participation was voluntary, users could exit at any time, and data were used only to improve the simulator and for academic research (not advertising). The main application unlocked only after clicking Agree. The tool was a policy and livelihood decision simulation for public education, not a clinical, diagnostic, or therapeutic intervention.

After consent, participants entered household data, completed the pre-survey when applicable, played through five simulation turns while viewing live indicators and logs, optionally requested AI debrief, and then completed the identity-tailored post-survey.

\printbibliography

\end{document}
